% chapters/08/praktikum.tex
\begin{praktikumbox}
Praktikum ini mengikuti Tugas Praktikum 06 pada RPS: NPC mengejar pemain di
ruang terbuka, dengan garis arah yang divisualisasikan agar keputusannya dapat
diamati.

\begin{langkah}
  \item \textbf{Jalankan programnya.} Ketik ulang \code{ch08_chase.py} dari
        Daftar~\ref{lst:ch08-chase} ke dalam folder \code{code}. Jalankan, lalu
        gerakkan pemain dengan WASD. Amati perubahan warna NPC dan garis arah
        yang menunjuk ke pemain.

  \item \textbf{Amati batas persepsi.} Berjalanlah menjauh sampai seluruh NPC
        berhenti, lalu dekati kembali. Catat jarak ketika NPC mulai bergerak,
        dan bandingkan dengan nilai jari-jari yang dipakai.

  \item \textbf{Ukur kecepatan pada berbagai jarak.} Jalankan program dengan
        masukan otomatis, lalu catat kecepatan terukur pada terminal.
\begin{kodekeluaran}
python ch08_chase.py --frames 400 --auto-run
\end{kodekeluaran}

  \item \textbf{Bandingkan dengan cara yang salah.} Jalankan kembali dengan
        pilihan \code{--naive}, lalu catat kecepatan terukurnya. Bandingkan
        kedua hasil dan jelaskan perbedaannya.
\begin{kodekeluaran}
python ch08_chase.py --frames 400 --auto-run --naive
\end{kodekeluaran}

  \item \textbf{Bandingkan jari-jari persepsi.} Jalankan program tiga kali
        dengan jari-jari 120, 240, dan 400 piksel. Catat kesan Anda mengenai
        tingkat kesulitan pada masing-masing nilai.

  \item \textbf{Amati kelemahan AI reaktif.} Tambahkan satu bongkahan tegak di
        tengah ruang permainan, lalu atur agar pemain dapat bersembunyi di
        baliknya. Perhatikan perilaku NPC selama pemain berada di balik
        bongkahan itu, dan catat pengamatan Anda.

  \item \textbf{Tulis catatan praktikum} yang memuat kedua tabel pengukuran,
        hasil pengamatan pada langkah keenam, dan kesimpulan mengenai kapan
        perilaku reaktif tidak lagi memadai.

  \item \textbf{Simpan pekerjaan} ke repositori Git dengan pesan penyimpanan
        yang menerangkan isinya.
\end{langkah}
\end{praktikumbox}

\begin{table}[htbp]
\centering
\caption{Lembar kerja Praktikum 06}
\label{tab:lembar-praktikum-06}
{\setlength{\tabcolsep}{4pt}\footnotesize%
\begin{tabularx}{\textwidth}{@{}C{4.0cm} C{3.2cm} C{3.2cm} Y@{}}
\toprule
\headrow \textbf{Cara dan pengukuran} & \textbf{Jarak saat diukur} & \textbf{Kecepatan terukur} & \textbf{Catatan} \\
\midrule
Arah dinormalkan & & & \\[1.4em]
\addlinespace
Vektor jarak dipakai langsung & & & \\[1.4em]
\addlinespace
\multicolumn{4}{@{}p{\textwidth}@{}}{\footnotesize\itshape
Jarak ketika NPC mulai bergerak (langkah 2): \dotfill} \\[1.4em]
\addlinespace
\multicolumn{4}{@{}p{\textwidth}@{}}{\footnotesize\itshape
Kesan tingkat kesulitan pada jari-jari 120, 240, dan 400 piksel: \dotfill}
\\[1.4em]
\addlinespace
\multicolumn{4}{@{}p{\textwidth}@{}}{\footnotesize\itshape
Perilaku NPC ketika pemain bersembunyi di balik bongkahan: \dotfill} \\[2.0em]
\bottomrule
\end{tabularx}}
\end{table}

\begin{catatabox}
Langkah keenam adalah langkah yang paling penting bagi bab berikutnya. Gejala
yang Anda catat di sana---NPC yang menekan bongkahan tanpa berusaha
mengelilinginya---merupakan titik tolak mesin keadaan pada Bab~IX dan pencarian
jalur pada Bab~X.
\end{catatabox}
